\ProvidesFile{bthesis-style-thesis-scu-master.tex}
  [2026/11/11 Thesis style: thesis-scu-master]

\RequirePackage{geometry}
\RequirePackage{fontspec}
\RequirePackage{xeCJK}
\RequirePackage{graphicx}
\RequirePackage{array}
\RequirePackage{booktabs}
\RequirePackage{caption}
\RequirePackage{titletoc}
\RequirePackage{fancyhdr}
\RequirePackage{indentfirst}
\RequirePackage{amsmath}
\RequirePackage{amssymb}
\RequirePackage{tikz}
\RequirePackage{hyperref}
\RequirePackage{etoolbox}

% 附件5的16开纸约197×273 mm；附件4示意图的版心为145×225 mm。
\geometry{papersize={197mm,273mm},textwidth=145mm,textheight=225mm,
  centering,headheight=14pt,headsep=5mm,footskip=10mm}
% 使用TeX Live自带字体，独立工程与Overleaf不依赖本机商业中文字体。
\setmainfont{texgyretermes-regular.otf}[BoldFont=texgyretermes-bold.otf,
  ItalicFont=texgyretermes-italic.otf,BoldItalicFont=texgyretermes-bolditalic.otf]
% Fandol不含GSUB，使用fontspec默认脚本而非请求不存在的CJK脚本。
\setCJKmainfont[Script=CustomDefault,BoldFont=FandolSong-Bold.otf,ItalicFont=FandolKai-Regular.otf]{FandolSong-Regular.otf}
\setCJKsansfont[Script=CustomDefault]{FandolHei-Regular.otf}
\setCJKfamilyfont{scusong}[Script=CustomDefault]{FandolSong-Regular.otf}
\setCJKfamilyfont{scuhei}[Script=CustomDefault]{FandolHei-Regular.otf}
\setCJKfamilyfont{scukai}[Script=CustomDefault]{FandolKai-Regular.otf}
\newcommand{\scuSong}{\CJKfamily{scusong}}
\newcommand{\scuHei}{\CJKfamily{scuhei}}
\newcommand{\scuKai}{\CJKfamily{scukai}}
\newcommand{\scuBodyFont}{\scuSong\fontsize{12bp}{20bp}\selectfont}
\newcommand{\scuSmallFont}{\scuSong\fontsize{10.5bp}{16bp}\selectfont}
\linespread{1}
\AtBeginDocument{\scuBodyFont}
\setlength{\parindent}{2em}
\setlength{\parskip}{0pt}
\xeCJKsetup{PunctStyle=plain}
\hypersetup{hidelinks,hypertexnames=false,bookmarksnumbered=true}

% 四级标题按issue要求顶格；紧邻标题不额外空行。
\ctexset{
  chapter={name={第,章},number=\chinese{chapter},format=\raggedright\scuHei\zihao{-3},
    beforeskip=20bp,afterskip=6bp,aftername=\quad,fixskip=true},
  section={name={第,节},number=\chinese{section},format=\raggedright\scuHei\zihao{4},
    beforeskip=20bp,afterskip=0bp,aftername=\quad},
  subsection={number=\chinese{subsection}、,format=\raggedright\scuHei\zihao{-4},
    beforeskip=20bp,afterskip=0bp,aftername=\hspace{0pt}},
  subsubsection={number=（\chinese{subsubsection}）,format=\raggedright\scuKai\zihao{-4},
    beforeskip=20bp,afterskip=0bp,aftername=\hspace{0pt}}
}
\setcounter{secnumdepth}{3}
\setcounter{tocdepth}{2}
\renewcommand{\contentsname}{目\quad 录}
\titlecontents{chapter}[0pt]{\scuBodyFont}{\thecontentslabel\quad}{}{\titlerule*[0.6pc]{.}\contentspage}
\titlecontents{section}[2em]{\scuBodyFont}{\thecontentslabel\quad}{}{\titlerule*[0.6pc]{.}\contentspage}
\titlecontents{subsection}[4em]{\scuBodyFont}{\thecontentslabel}{}{\titlerule*[0.6pc]{.}\contentspage}
\fancypagestyle{scu-main}{%
  \fancyhf{}%
  \fancyhead[C]{\scuSong\fontsize{9bp}{12bp}\selectfont 四川大学硕士学位论文}%
  \fancyfoot[C]{\scuSmallFont\thepage}%
  \renewcommand{\headrulewidth}{0.4pt}%
  \renewcommand{\footrulewidth}{0pt}%
}
\fancypagestyle{plain}[scu-main]{}
\DeclareCaptionFont{scucaption}{\scuHei\fontsize{10.5bp}{16bp}\selectfont}
\captionsetup{font=scucaption,labelsep=quad,skip=10bp}
\setlength{\textfloatsep}{20bp}
\setlength{\intextsep}{20bp}
\AtBeginEnvironment{tabular}{\scuSmallFont}
\AtBeginEnvironment{tabular*}{\scuSmallFont}

% 附件6：连续阿拉伯数字上标脚注，小五宋体；跨章不重置。
\makeatletter
% 章节、列表和浮动体结束时会重新调用normalsize，须同时固定其基线。
\renewcommand{\normalsize}{\@setfontsize\normalsize{12bp}{20bp}%
  \abovedisplayskip=10bp\belowdisplayskip=10bp%
  \abovedisplayshortskip=0bp\belowdisplayshortskip=6bp}
\@removefromreset{footnote}{chapter}
\renewcommand{\footnotesize}{\scuSong\fontsize{9bp}{12bp}\selectfont}
\makeatother
\renewcommand{\thefootnote}{\arabic{footnote}}
% thebibliography自行生成标题，调用方不再另写chapter*。
\AtBeginEnvironment{thebibliography}{\scuSmallFont}
\apptocmd{\thebibliography}{\addcontentsline{toc}{chapter}{参考文献}\scuSmallFont}{}{}
\newcommand{\scuReferenceGroup}[1]{\par\addvspace{10bp}\item[] {\scuHei #1}\par\nopagebreak}

\providecommand{\scuTitleZh}{硕士学位论文题目}
\providecommand{\scuTitleEn}{Title of the Master's Thesis}
\providecommand{\scuAuthor}{作者姓名}
\providecommand{\scuSupervisor}{导师姓名及职称}
\providecommand{\scuSchool}{四川大学高分子科学与工程学院}
\providecommand{\scuMajor}{材料科学与工程}
\providecommand{\scuDirection}{高分子材料}
\providecommand{\scuCompletionDate}{二〇二六年十一月}
\providecommand{\scuDegreeDate}{填写授予学位日期}
\providecommand{\scuKeywordsZh}{高分子；链构象；尺度关系}
\providecommand{\scuKeywordsEn}{polymer; chain conformation; scaling}
\providecommand{\scuWordmark}{assets/branding/scu-wordmark.png}

\newcommand{\scuMakeCover}{%
  \thispagestyle{empty}%
  \begin{center}
    \vspace*{8mm}
    \includegraphics[width=76mm]{\scuWordmark}\par
    \vspace{12mm}
    {\scuSong\zihao{1} 硕\quad 士\quad 学\quad 位\quad 论\quad 文\par}
    \vspace{18mm}
    {\scuHei\zihao{3} \scuTitleZh\par}
    \vspace{12mm}
    {\scuBodyFont 作者：\scuAuthor\quad 完成日期：\scuCompletionDate\par}
    \vspace{14mm}
    \renewcommand{\arraystretch}{1.8}
    \begin{tabular}{rl}
      培养单位 & \scuSchool\\
      指导教师 & \scuSupervisor\\
      专业 & \scuMajor\\
      研究方向 & \scuDirection\\
      授予学位日期 & \scuDegreeDate
    \end{tabular}
  \end{center}
  \clearpage
}
\newcommand{\scuAbstractZhHeading}{%
  \phantomsection\addcontentsline{toc}{chapter}{中文摘要}%
  \begin{center}
    \vspace*{40bp}
    {\scuSong\zihao{-2}\scuTitleZh\par}
    \vspace{20bp}
    {\scuBodyFont\scuMajor\par}
    \vspace{20bp}
    {\scuHei\zihao{-4}研究生：{\scuKai\scuAuthor}\quad 指导教师：{\scuKai\scuSupervisor}\par}
    \vspace{40bp}
    {\scuHei\zihao{4}摘\quad 要\par}
  \end{center}\vspace{20bp}
}
\newcommand{\scuAbstractEnHeading}{%
  \phantomsection\addcontentsline{toc}{chapter}{Abstract}%
  \begin{center}
    \vspace*{40bp}
    {\bfseries\zihao{-2}\scuTitleEn\par}
    \vspace{40bp}
    {\bfseries\zihao{4}ABSTRACT\par}
  \end{center}\vspace{20bp}
}
